\documentclass[11pt]{article}

\usepackage[T1]{fontenc}
\usepackage{lmodern}
\usepackage[a4paper,margin=1in]{geometry}
\usepackage{amsmath,amssymb}
\usepackage{booktabs}
\usepackage{array}
\usepackage{enumitem}
\usepackage{underscore}
\usepackage{microtype}
\usepackage[dvipsnames]{xcolor}
\usepackage[colorlinks=true,linkcolor=black,urlcolor=NavyBlue,citecolor=NavyBlue]{hyperref}

\setcounter{secnumdepth}{3}
\setcounter{tocdepth}{2}

\newcommand{\code}[1]{\texttt{\small #1}}
\newcommand{\pending}{\textit{[Draft pending.]}}
\newcommand{\field}[1]{\paragraph{#1.}}

\setlength{\parskip}{0.4em}
\setlength{\parindent}{0pt}

\title{\textbf{Behavioral responses in the top-tax project}\\[0.3em]
\large A methodology}
\author{Budget Lab at Yale \\ \small Tax-Simulator}
\date{Draft, \today}

\begin{document}
\maketitle

\begin{abstract}
\noindent This paper documents the behavioral responses and cross-base
interactions used in the top-tax project, which estimates the additional
revenue available from the top of the income and wealth distribution. It
describes each channel under a common template---mechanism, parameters,
calibration, interactions, and limitations---and the architecture that lets the
channels interact on shared record-level state. This is a working draft: the
capital-gains realization channel (Section~\ref{sec:realization} and
Appendix~\ref{sec:app-realization}) is complete and serves as the template
exemplar; the remaining channels are drafted or outlined to the same structure.
\end{abstract}

\tableofcontents
\bigskip

%======================================================================
\section{Introduction}
\label{sec:intro}

Several current proposals draw revenue from the same taxpayers through distinct
statutes: top ordinary rates, capital-gains rates, the treatment of gains at
death, the estate tax, a wealth tax, the corporate tax, and repeal of the
qualified business income deduction. Scoring each in isolation and summing
overstates their combined yield, for two reasons.

First, the bases overlap mechanically. A given stock of wealth is capital income
during life, an estate at death, a wealth-tax base in the interim, and corporate
profit upstream. A dollar collected under one statute is unavailable to another,
so summing standalone estimates double-counts it.

Second, taxpayers respond. They retime realizations, convert ordinary income
into deferred gains, shift income across the corporate and pass-through
boundary, conceal assets, and finance tax payments out of saving. These
responses move income between bases and out of the taxable total, and a response
induced by one instrument alters the base on which the others operate.

The model estimates these effects on record-level balance sheets rather than
through an assumed aggregate elasticity. Each channel described below is a
structural response measured against the microdata; the object of interest is
their interaction (Section~\ref{sec:architecture}), which is why the project
reports the conventional estimate as a share of the static estimate and a
by-base destination ledger rather than a single headline elasticity.

Two of the channels documented here, the wealth-financing channel
(Section~\ref{sec:bathtub}) and corporate incidence
(Section~\ref{sec:corp}), are mechanical rather than elasticity-driven. They are
included because they transmit a policy's effect from one base to another, which
is the interaction the project measures.

%======================================================================
\section{Model structure}
\label{sec:structure}

\subsection{Static and conventional passes}
Each non-baseline scenario is run twice. The \emph{static} pass holds taxpayer
inputs at baseline levels and recomputes tax under the reform; it is the
law-only counterfactual. The \emph{conventional} pass activates the behavioral
channels. The reported quantities are differences between the two: the static
estimate is the mechanical intent, and the conventional estimate is the amount
that survives behavior. Interactions are reported through receipts as static
minus conventional deltas. Distribution tables are static-sourced by convention.

\subsection{Two channel types}
\label{sec:types}
The model implements behavior in two forms.

\begin{table}[h]
\centering
\small
\begin{tabular}{@{}p{0.22\textwidth}p{0.70\textwidth}@{}}
\toprule
\textbf{Type} & \textbf{Implementation} \\
\midrule
Behavior module &
An R file at \code{config/scenarios/behavior/\{var\}/\{name\}.R} exposing a
\code{do_\{var\}()} hook, dispatched on the behavioral pass. It modifies
tax-unit inputs; tax is then recomputed. Examples: realization, income
conversion, entity shifting, evasion, wealth avoidance, charity, employment. \\
\addlinespace
Mechanical applier &
A fixed step invoked at the head of the conventional pass in
\code{src/sim/run.R}, with no hook and no elasticity dispatch. Examples: the
wealth-financing channel and corporate incidence. \\
\bottomrule
\end{tabular}
\end{table}

\subsection{Run order}
\label{sec:order}
Because the channels touch overlapping bases, their order is part of the
specification. The conventional sequence is
\[
\text{corporate incidence} \rightarrow \text{wealth haircut} \rightarrow
\underbrace{\text{realization} \rightarrow \sigma \rightarrow
\text{entity shifting} \rightarrow \text{evasion}}_{\text{fixed module order}}
\rightarrow \texttt{do\_taxes}.
\]
The module order is asserted at runtime so that each dollar is adjusted at most
once: realizations are set before conversion prices them, conversion before the
entity margin re-routes the remainder, and the reporting channels (evasion,
avoidance) run last, after which only \code{do_taxes} reads the frame
(Section~\ref{sec:firewall}).

%======================================================================
\section{Interaction architecture}
\label{sec:architecture}

Three structures let the channels interact on common state rather than through
hand-set cross-elasticities.

\subsection{Shared objects}
\label{sec:shared}
Several channels are priced from the same internal quantities:
\begin{itemize}[leftmargin=1.4em]
\item $\tau_{eq}(a,t)$, the present value of tax per dollar entering the
unrealized-gain state, produced by the realization machinery
(Section~\ref{sec:realization}). It prices the conversion exit
(Section~\ref{sec:conversion}) and, on the roadmap, the entity-shifting deferral
value (Section~\ref{sec:entity}).
\item $F$, the death-forgiveness value of an unrealized gain, which depends on
the realization-at-death regime and on estate exposure
(Section~\ref{sec:estate-offset}). It links realization to the estate tax.
\item $h=\tau_w\,\tau_{cg}$, the wealth-tax carry on a deferred gain
(Section~\ref{sec:carry}), which couples a wealth tax into the realize-or-hold
decision.
\end{itemize}
Because these objects are shared, a change to one instrument propagates to the
others through the model's own state.

\subsection{The reporting--real separation}
\label{sec:firewall}
The governing rule is that cash is reported and balance sheets are real. A
reporting response (evasion, wealth avoidance, estate concealment) affects only
the cash tax paid, and through it the financing flow of the wealth channel
(Section~\ref{sec:bathtub}). It does not affect the economic balance sheet
(\code{value.*}), the gain stock, cohort assignment, heir inheritances, or
corporate exposure. A real response (corporate markdown, wealth haircut,
realization, conversion, entity shifting) flows into every base, including
reported ones.

The architecture enforces this. Every real-side consumer reads raw input data,
static detail, or a preserved raw column, never the behavior-modified
conventional frame. Reporting modules write only to tax-computation inputs and
to documented isolation points (the materialized \code{net_worth} column, the
\code{estate_concealed_frac} input to \code{calc_estate}). This is why the
reporting modules are ordered last (Section~\ref{sec:order}). Each channel below
is marked \textsc{real} or \textsc{reporting}.

\subsection{Composition of the taxable-income elasticity}
\label{sec:eti}
The top-bracket elasticity of taxable income (ETI) is composed from three
modules---income conversion (Section~\ref{sec:conversion}), entity shifting
(Section~\ref{sec:entity}), and evasion (Section~\ref{sec:evasion})---together
with charity. One parameter, $\sigma$, is calibrated as a residual so that the
modeled top-ordinary ETI matches an external target of approximately $0.25$,
conditional on the rest of the stack.

A consequence is that any change to a margin inside the ordinary-rate ETI bundle
requires $\sigma$ to be re-derived. Margins outside the bundle---the wealth and
estate terms in the realization problem, the estate own-rate response, and
capital-gains realization itself---do not affect $\sigma$. This distinction
governs calibration maintenance and is monitored by an automated staleness check
(Section~\ref{sec:limitations}).

%======================================================================
\section{Behavioral channels}
\label{sec:channels}

Each channel is documented under seven headings:
\begin{enumerate}[leftmargin=1.6em,itemsep=0.1em]
\item \emph{What it captures} --- the economic response, stated briefly.
\item \emph{Mechanism} --- how the model produces it.
\item \emph{Type and location} --- module or mechanical; \textsc{real} or
\textsc{reporting}; code path.
\item \emph{Parameters} --- the pinned values and override switches.
\item \emph{Calibration and evidence} --- how the parameters are set and the
literature anchor.
\item \emph{Interactions} --- the other channels and bases it touches.
\item \emph{Limitations} --- what it omits.
\end{enumerate}
Medium-depth content appears here; full derivations are in
Appendix~\ref{sec:appendix}.

%----------------------------------------------------------------------
\subsection{Realization and timing at death}

\subsubsection{Capital-gains realization: the entropy Bellman}
\label{sec:realization}

\field{What it captures}
The share of a mechanical capital-gains rate increase that survives once
realizations respond. Investors defer gains when rates rise, and the value of
deferral depends on the treatment of the gain at death. The channel models both
the long-run level of realizations and the short-run retiming around a rate
change.

\field{Mechanism}
Investors are grouped into representative cells defined by age cohort and
within-age net-worth percentile. For each cell the model solves a
realize-or-hold problem over the gain stock under an entropy realization cost.
Each dollar of unrealized gain carries a non-tax motive to sell---liquidity,
rebalancing, or consumption---drawn from an exponential reservation
distribution; the holder realizes the dollar when that motive exceeds the tax
wedge between realizing now and deferring. The realization rate is the survival
function of that distribution. The first-order condition is
\begin{equation}
\label{eq:foc}
r_D^{\,j} = r_D^{\,B}\exp\!\big(-\eta\,(MC^{\,j}-MC^{\,B})\big),
\qquad \text{clipped to } [0,1],
\end{equation}
where the marginal cost of holding rather than realizing is
\begin{equation}
\label{eq:mc}
MC = \tau + \beta(1-m)\,W_{\text{next}} + \beta m F,
\qquad F = (1-c_\phi)\,\tau.
\end{equation}
Here $\tau$ is the current-year rate, $W_{\text{next}}$ the discounted
continuation value of an unrealized dollar under survival probability $1-m$, and
$F$ the discounted death-forgiveness value under mortality $m$. The
realization-at-death regime enters through $c_\phi$, the share of the gain's
tax burden the holder internalizes at death: $c_\phi=0$ under step-up ($F=\tau$),
$c_\phi=1$ under deemed realization ($F=0$), and an intermediate value under
carryover.

Since $\mathrm{d}\ln r_D/\mathrm{d}MC = -\eta$ across the whole pool, $\eta$ is
the long-run realization semi-elasticity directly, with no inert floor or
responsive-share deflation; this is the single-pool property of the current
specification (Appendix~\ref{sec:app-realization}). A separate short-run overlay
lets a calibrated fraction of realizations retime by one year toward the
lower-wedge year. The overlay nets to zero under a uniform permanent shock and
therefore leaves the long-run response unchanged.

\field{Type and location}
Behavior module. The applier \code{do_kg_dynamics}
(\code{config/scenarios/behavior/kg_dynamics/turnover.R}) is an allocator that
reads per-year cell state and translates it to per-record \code{kg_lt}
adjustments. State is produced by a pre-pass (\code{kg_dyn_run_bathtub_pass}) in
\code{src/sim/kg_dynamics.R}; shared cohort primitives are in
\code{src/sim/cohort_bathtub.R}. \textsc{real}.

\field{Parameters}
\begin{table}[h]
\centering\small
\begin{tabular}{@{}llp{0.42\textwidth}l@{}}
\toprule
Parameter & Value & Meaning & Override \\
\midrule
$\eta$ & $2.4825$ & long-run realization semi-elasticity & \code{KG_ETA} \\
timeable share & $0.2542$ & fraction of realizations that retimes by one year & \code{KG_TIMEABLE_SHARE} \\
deemed-avoidance & $0.25$ & valuation discount on deemed realizations (data, not law) & \code{KG_DEEMED_AVOIDANCE} \\
$c_\phi$ & $0/\theta/1$ & regime death-burden share & --- \\
spec version & $3$ & single pool & --- \\
\bottomrule
\end{tabular}
\end{table}

\field{Calibration and evidence}
$\eta$ is pinned on the full simulator. The calibration harness measures the
realized elasticity $E_{\text{full}}(\eta)=\mathrm{d}\log R/\mathrm{d}\tau_{rw}$
at simulation year 30 across a grid of $\eta$ values; the surface is linear
through the origin, so a few points determine it. Inverting the fitted line for
the target $E_{\text{full}}^{\star}=-0.6/0.238=-2.52$ (a top-rate realization
elasticity over the top-rate divisor, following the JCT and CBO convention and
the estimates of Dowd, McClelland, and Muthitacharoen) at slope $1.0155$ gives
$\eta^{\star}=2.4825$. The fitted coefficients are recorded so that the next
re-pin is arithmetic. The timeable share is calibrated separately against the
short-run announcement moment given $\eta$. The full procedure is in
Appendix~\ref{sec:app-realization}.

\field{Interactions}
The channel produces $\tau_{eq}$ and $F$ (Section~\ref{sec:shared}), which price
income conversion (Section~\ref{sec:conversion}) and entity shifting
(Section~\ref{sec:entity}). The death regime links realization to the estate
tax: closing step-up raises realizations. The wealth-carry term
(Section~\ref{sec:carry}) and the estate death-value offset
(Section~\ref{sec:estate-offset}) enter $MC$ and $F$, coupling the wealth and
estate instruments into the realization decision.

\field{Limitations}
The five tracked wealth classes are collapsed into one asset bucket;
per-asset-class disaggregation is on the roadmap. The revenue-maximizing grid
argmaxes at the current boundary (a statutory peak near 44--45\%) and should be
extended to identify the turning point. The deemed-avoidance discount is a
data calibration consistent with JCT scoring, not a behavioral response.

\subsubsection{Wealth-carry}
\label{sec:carry}

\field{What it captures}
A wealth tax raises the cost of deferring a gain, because the unpaid
capital-gains tax on an unrealized gain remains in the wealth-tax base and is
taxed at $\tau_w$ each year it is deferred. This unlocks realizations in
wealth-taxed cells that would otherwise defer, and it lowers the value of the
conversion deferral path (Section~\ref{sec:conversion}).

\field{Mechanism}
A per-year carrying cost $h=\tau_w\,\tau_{cg}$ is added to the marginal cost of
holding in~\eqref{eq:mc} and to the $\tau_{eq}$ recursion. It is computed at the
record level and then aggregated, not as a product of separately averaged rates,
because $\tau_w$ and $\tau_{cg}$ are positively correlated at the top. The term
applies on the survival branch only; in-year realizers and decedents pay no
carry that year.

\field{Type and location}
A structural term inside the realization module (\code{src/sim/kg_dynamics.R}).
\textsc{real}.

\field{Parameters}
None free. $h$ is the product of two marginal rates already in the model,
$\tau_w$ from \code{mtr_net_worth} and $\tau_{cg}$ from \code{mtr_kg_lt}. Under
current law $\tau_w=0$ for essentially all cells, so $h=0$ is an exact no-op,
which is why the term did not change the specification version or the $\eta$
calibration.

\field{Calibration and evidence}
No separate calibration. The magnitude follows from the deferral arithmetic: the
per-year benefit of deferral is approximately $r_{\text{real}}\,\tau$, so a
one-percent wealth tax offsets roughly half of it at $r_{\text{real}}\approx2\%$.

\field{Interactions}
Couples the wealth-tax instrument into realization and into the conversion price.

\field{Limitations}
The response is aggregated at the age-cell level rather than the record level,
which understates the record-level response by Jensen's inequality (conservative,
never an overstatement); the measured understatement ranges from about 8\% to
33\% depending on the parameter corner. An exposed/unexposed split closing the
gap was considered and not adopted. The understatement should be noted when
reporting wealth and capital-gains corners jointly.

\subsubsection{Estate death-value offset}
\label{sec:estate-offset}

\field{What it captures}
For an estate-taxable decedent, capital-gains or deemed tax paid at or before
death reduces the taxable estate, so the effective rate near death is
approximately $\tau_{cg}(1-\tau_e)$. Including this stops the model from
overstating lock-in for estate-exposed cells when the deemed and estate switches
are active. \pending

\field{Mechanism, parameters, and limitations}
Gain-weighted cell aggregation of a switch-gated marginal estate rate entering
$F$ and $MC$; baseline death value $F_B=\tau_B(1-e_B)$, paired by leg. Inside
the realization module; \textsc{real}. \pending

%----------------------------------------------------------------------
\subsection{Taxable-income base-shifting}

\subsubsection{Income conversion ($\sigma$)}
\label{sec:conversion}

\field{What it captures}
Conversion of labor and ordinary income into the deferred unrealized-gain state
(the founder-equity and carried-interest path), a component of the top-ordinary
ETI. \pending

\field{Summary}
Module \code{do_conversion} (\code{config/scenarios/behavior/conversion/sigma.R};
shared function \code{src/sim/sigma_conversion.R}); \textsc{real}. Parameter
$\sigma=0.16$ (\code{SIGMA_CONV}), re-derived in 2026, superseding an earlier
$0.08$. Calibrated as a residual to a top-ordinary ETI of $0.25$
(Section~\ref{sec:eti}). Consumes $\tau_{eq}$; a change to the entity, evasion,
or charity margins requires re-derivation. The conversion path is currently
priced free of the corporate layer. \pending

\subsubsection{Entity shifting}
\label{sec:entity}

\field{What it captures}
Business income relocating across the corporate and pass-through boundary in
response to the corporate-to-individual rate differential, following Pearce and
Prisinzano. \pending

\field{Summary}
Module \code{do_entity_shifting}
(\code{config/scenarios/behavior/entity_shifting/pearce_prisinzano.R});
\textsc{real}. The shifting semi-elasticity is presently a fixed value of
$0.25$. This is a known limitation: it should be derived from the model's own
deferral value $\tau_{eq}/\tau$, which would make the corporate retention shelter
respond to the realization-at-death regime rather than remaining invariant to
it. \pending

\subsubsection{Evasion}
\label{sec:evasion}

\field{What it captures}
Under-reporting of self-employment, partnership, S-corporation, and rental
income that rises with the net-of-tax rate; the leakage component of the ETI.
\pending

\field{Summary}
Module \code{do_evasion} (\code{config/scenarios/behavior/evasion/debacker.R});
\textsc{reporting}. Net-of-tax elasticities: Schedule C/F $0.046$, partnership
and S-corporation $0.052$, rent $0.040$. Wages, interest, and dividends do not
respond, reflecting information reporting. Evidence from DeBacker, Heim, and
Yuskavage. Feeds the hidden-ledger link (Section~\ref{sec:avoidance}).
Limitations: positive legs only; the overstated-loss margin and a top-graded
multiplier are deferred. \pending

%----------------------------------------------------------------------
\subsection{Wealth and estate own-rate responses}

\subsubsection{Wealth avoidance and the hidden ledger}
\label{sec:avoidance}

\field{What it captures}
Reduction in reported taxable wealth under a wealth tax, through legal valuation
management and concealment, with concealment made consistent across bases so
that a concealed dollar also escapes income, capital-gains, and estate tax.
\pending

\field{Summary}
Module \code{do_wealth} (\code{config/scenarios/behavior/wealth/avoidance.R});
\textsc{reporting}, writing only the \code{net_worth} isolation point and
\code{estate_concealed_frac}. Semi-elasticities: marketable wealth $-7$,
closely-held wealth $-17$. Concealment shares: public $1.0$, private $0.5$.
Migration is subsumed here as a ceiling (Section~\ref{sec:other}). \pending

\subsubsection{Estate own-rate response}
\label{sec:estate-ownrate}

\field{What it captures}
Reduction in reported gross estate in response to a change in the
net-of-estate-tax rate. \pending

\field{Summary}
Part of \code{do_wealth}; \textsc{reporting}, via \code{estate_concealed_frac}.
Elasticity $\varepsilon=0.16$ (\code{ESTATE_REPORT_EPS}; Kopczuk--Slemrod band
$0.10$--$0.22$), in the net-of-tax form
$\text{retained}=\big((1-\tau_S)/(1-\tau_B)\big)^{\varepsilon}$, keyed to the
change in the marginal estate rate rather than its level. Evidence from Kopczuk
and Slemrod. A rate-responsive charitable-bequest margin is not yet built.
\pending

%----------------------------------------------------------------------
\subsection{Saving and financing}

\subsubsection{The wealth-financing channel}
\label{sec:bathtub}

\field{What it captures}
A share $s=1-\text{MPC}$ of the net above-baseline during-life tax is financed
out of wealth rather than consumption. It compounds over time and enters the
estate and capital-income bases at death. This is the channel through which a
during-life tax reduces future wealth-based revenue. \pending

\field{Mechanism}
A cohort recurrence with kernel and forcing
\[
G(a,p,t) = \big(1+r_{\text{total}}(t)\big) - s\,(\tau y + \tau_w),
\qquad
F = \Delta T^{0} - \Delta Y_{\text{exog}}.
\]
\pending

\field{Summary}
Mechanical applier (\code{src/sim/wealth_dynamics.R},
\code{src/sim/cohort_bathtub.R}), run at the head of the conventional pass;
\textsc{real}. Parameters: the saving-share profile $s(\text{age},\text{net-worth
percentile})$, calibrated from about $0.1$ at the bottom to $0.80$ at the top; a
transition matrix; a cap $\code{fmax}=0.9$; and $r_{\text{total}}$ equal to
nominal GDP-per-capita growth. It is a financing response, not a rate-of-return
saving elasticity. It requires a dedicated pass, roughly doubling compute.
\pending

%----------------------------------------------------------------------
\subsection{Corporate incidence and labor-market margins}

\subsubsection{Corporate incidence}
\label{sec:corp}

\field{What it captures}
The record-level incidence of a corporate rate change: cuts to dividend,
interest, rent, and pass-through flows; an equity markdown on exposed asset
stocks; a capital-gains adjustment; dissaving through the wealth channel; and an
endogenous individual-tax offset. \pending

\field{Summary}
Mechanical, revenue-side applier (\code{src/sim/corp_incidence.R}) with
fail-closed activation on \code{corporate_meta.yaml}; \textsc{real}. Parameters
are hardcoded \code{CORP_*} constants, several of which are placeholders pending
measurement and disclosed as such. The endogenous offset makes the corporate row
sensitive to stacking order. The most complete existing documentation is
\code{other/corporate_incidence/CONSIDERATIONS.md} and \code{FORMAL_MODEL.md}.
\pending

\subsubsection{Charity, employment, and migration}
\label{sec:other}

\begin{itemize}[leftmargin=1.4em]
\item \textbf{Charity.} Module \code{do_charity}; tax-price elasticity $-0.5$ on
cash giving. The appreciated-asset margin, through which a higher capital-gains
rate raises the subsidy to donating securities, is not built. \pending
\item \textbf{Employment.} Module \code{do_employment}
(\code{config/scenarios/behavior/employment/bastian.R}); extensive-margin wage
elasticities following Bastian. Second-order for the top-tax question; documented
for completeness. \pending
\item \textbf{Migration.} No standalone channel. It is subsumed as a ceiling
within the wealth-avoidance semi-elasticity and disclosed as such.
\item \textbf{Excluded by convention.} Real labor supply and a rate-of-return
saving elasticity are excluded; the wealth channel is financing, not saving.
\end{itemize}

%======================================================================
\section{Limitations and open items}
\label{sec:limitations}

Two facts frame this section. First, the realization problem prices from a single
tax rate plus the shared $F$, $h$, and estate terms, so any mispriced cross-rate
content appears as a missing interaction rather than as noise. Second, $\sigma$
is a residual, so a gap inside the ETI bundle carries a re-calibration cost that
a gap outside it does not (Section~\ref{sec:eti}).

\begin{table}[h]
\centering\small
\begin{tabular}{@{}p{0.30\textwidth}cp{0.34\textwidth}l@{}}
\toprule
Item & Ref. & Effect & Status \\
\midrule
Entity-shifting fixed at $0.25$ & \ref{sec:entity} & shelter invariant to death regime & open, in-bundle \\
Corporate constants placeholder & \ref{sec:corp} & corporate estimate and ETR band & disclosed \\
Wealth-carry cell aggregation & \ref{sec:carry} & understates by 8--33\%, conservative & accepted \\
Appreciated-asset giving margin & \ref{sec:other} & understates giving and lock-in at high rates & open, in-bundle \\
Migration not explicit & \ref{sec:other} & subsumed as a ceiling & by convention \\
Corporate layer absent from $\sigma$ & \ref{sec:conversion} & conversion priced corporate-free & open, in-bundle \\
Estate own-rate coverage & \ref{sec:estate-ownrate} & no charitable-bequest response & partial \\
Evasion coverage & \ref{sec:evasion} & positive legs only; flat top multiplier & partial \\
\bottomrule
\end{tabular}
\end{table}

\paragraph{Calibration maintenance.}
The residual parameters, $\sigma$ and $\eta$ under its conditioning assumptions,
are monitored by an automated check: a low-cost internal-moment drift test on
each run, and a $\sigma$ re-derivation triggered by any commit that touches the
conditioning set (entity shifting, evasion, the charity elasticity, the
conversion pool, or the input-data vintage). Reference values are stored in
\code{other/kg_model_tests/calibration_reference.csv}. The check detects drift;
it does not re-pin automatically.

%======================================================================
\section{References}
\label{sec:references}
\emph{To be completed.} Kopczuk and Slemrod (2001); Dowd, McClelland, and
Muthitacharoen (2015); DeBacker, Heim, and Yuskavage; Pearce and Prisinzano
(2018); Bastian (2023); Saez and Zucman on wealth-tax avoidance; and the JCT and
CBO realization-elasticity conventions.

%======================================================================
\appendix
\section{Derivations and calibration}
\label{sec:appendix}

\subsection{Capital-gains realization}
\label{sec:app-realization}

\paragraph{The problem.}
For a representative cell, let $r_D$ be the discretionary realization rate on the
gain pool. Each unrealized dollar carries a non-tax reservation benefit $b\ge0$
to selling, drawn from an exponential distribution, and the holder realizes when
$b$ exceeds the tax wedge between realizing now and holding. The entropy
realization cost $C(r_D)$ is the cost whose marginal is the inverse survival
function of that exponential,
\begin{equation}
\label{eq:cprime}
C'(r_D) = \tfrac{1}{\eta}\ln\!\big(r_D/r_D^{\,B}\big),
\qquad C'(r_D^{\,B}) = 0,
\end{equation}
so the baseline rate is the zero-cost reference. The holder equates the net
benefit of realizing, $\kappa-MC$, to $C'(r_D)$, with $\kappa$ the reservation
constant and $MC$ as in~\eqref{eq:mc}.

\paragraph{Two passes.}
Baseline inversion gives $C'(r_D^{\,B})=0$, hence $\kappa=MC^{B}$ exactly,
interior and at the corner. The scenario condition $\kappa-MC^{j}=C'(r_D)$ then
yields~\eqref{eq:foc}. Only the upper clip binds, and cells with $r_D^{\,B}=0$
stay at zero. Since $\mathrm{d}\ln r_D/\mathrm{d}MC=-\eta$ everywhere and the
whole pool responds, $\eta$ is the aggregate long-run semi-elasticity and the
revenue-maximizing rate is approximately $1/\eta$.

\paragraph{Single-pool specification.}
An earlier specification split baseline realizations into a responsive share and
an inert share, the latter a point mass of fully constrained sellers. That floor
was an artifact of an earlier quadratic cost with no tail, and it removed the
interior revenue maximum. The exponential reservation distribution already
represents the heterogeneity the floor stood in for, so the current
specification drops it: one pool, with a level margin from the Bellman and a
timing margin from the overlay.

\paragraph{Timing overlay.}
A fraction $f$ of baseline realizations retimes across a one-year window toward
the lower-wedge year, with move-share
$\operatorname{clamp}(\Delta\text{wedge}/\text{ref\_wedge},0,1)$ and
$\text{ref\_wedge}=0.05$. Composed as a net shift
$r_S=r_{\text{ord},S}+(r_{\text{plan},S}-r_{\text{plan},B})$, it nets to zero
under a uniform permanent shock, so $\eta$ is identified independently of $f$.

\paragraph{Calibration of $\eta$.}
On the full simulator, measure
$E_{\text{full}}(\eta)=\log(R_{\text{shock}}/R_{\text{base}})/\Delta\tau_{rw}$ at
simulation year 30 across a grid; the surface is linear through the origin.
Invert for $\eta^{\star}=|E_{\text{full}}^{\star}|/\text{slope}$ with
$E_{\text{full}}^{\star}=-2.52$ and slope $1.0155$, giving $\eta^{\star}=2.4825$
(grid $\eta\in\{2,2.3992,3\}$ mapping to
$E_{\text{full}}\in\{-2.05,-2.44,-3.03\}$). Provenance is guarded: a stale value
warns, and a strict-calibration flag halts.

\bigskip
\noindent\emph{Appendix derivations for the remaining channels follow at this
depth.}

\section{Parameter and provenance table}
\label{sec:app-params}

\begin{table}[h]
\centering\small
\begin{tabular}{@{}llll@{}}
\toprule
Constant & Value & Reference moment & Derived \\
\midrule
\code{KG_DYN_DEFAULT_ETA} & $2.4825$ & full-sim $E_{\text{full}}=-2.52$ / slope $1.0155$ & 2026-07-12 \\
\code{KG_DYN_TIMEABLE_SHARE} & $0.2542$ & full-sim short-run moment & 2026-07-09 \\
\code{SIGMA_CONV} & $0.16$ & top-ordinary ETI $\approx0.25$ & 2026-07-12 \\
\code{ESTATE_REPORT_EPS} & $0.16$ & Kopczuk--Slemrod (band $0.10$--$0.22$) & 2026-07-12 \\
Evasion (C/PT/rent) & $0.046/0.052/0.040$ & DeBacker--Heim--Yuskavage & --- \\
Wealth avoidance (mkt/closely-held) & $-7/-17$ & author-accepted & 2026-07-08 \\
Concealment share (public/private) & $1.0/0.5$ & consistency across bases & 2026-07-08 \\
Entity shifting & $0.25$ (stub) & should be $\tau_{eq}/\tau$ & --- \\
\bottomrule
\end{tabular}
\end{table}

\noindent The reference-moment column is also the source for the staleness check
(\code{calibration_reference.csv}).

\end{document}
